A measurement that says the interval is too narrow is only useful if something
uses it. The package ships the map, and the fit consults it.

After the optimizer returns, the fit looks its own regime up: the event count, the
fitted branching ratio, the smallest fitted decay rate and the observation window
place it in the measured grid, and the nearest cell returns the measured $q$ and
its Monte Carlo standard error for each parameter block. The decision then follows
one rule. If a block's two-standard-error interval around the measured $q$ excludes
one, the direction of the miscalibration is resolved at the precision the study
supports, and the fit replaces the Hessian interval with a parametric bootstrap. If
the map does not cover the regime, the Hessian interval is used and labelled
uncalibrated, because no measurement supports anything better. If the map reports
the estimator undefined, no interval is reported at all.

The switch has a cost, and the cost is bounded. The bootstrap refits the model 200
times, which at the largest measured configuration costs \AutoPoolCost{}~s on the
default process pool and \AutoSerialCostMin{}~min on one worker. A fit therefore
projects the cost before spending it, and refuses to exceed \AutoLimit{}~s: above
that ceiling it keeps the Hessian interval, says what was indicated and what it
would have cost, and warns. The projection is a power law fitted to the four
measured serial configurations of the capability table, multiplied by the measured
efficiency of the process pool, and it reproduces every measured row to within
three percent.

Whether the switch helps is a separate measurement, pre-registered before it ran:
\AutoCells{} cells at $N = 2000$, \AutoReplications{} replications each, both arms
fitted to the same simulated event stream so the comparison is paired.
\texttt{auto} chose the bootstrap in \AutoBootstrapShare{} percent of replications.

\begin{table}[t]
\centering
\caption{Coverage of nominal-95\% intervals under the Hessian interval and under
\texttt{auto}, with the paired difference $\delta = |\text{cov}_{\text{auto}} -
0.95| - |\text{cov}_{\text{hessian}} - 0.95|$. Negative $\delta$ means
\texttt{auto} is closer to nominal.}
\begin{tabular}{lrrr}
\toprule
Cell & coverage, Hessian & coverage, auto & difference \\
\midrule
\multicolumn{4}{c}{\emph{auto-v1 has not run}} \\
\bottomrule
\end{tabular}

\end{table}

H-A1, that \texttt{auto} lands closer to nominal coverage than the Hessian does, is
\AutoHoneLabel{}. At $\rho = 0.3028$ the pooled coverage moves from \AutoLowRhoCovH{}
to \AutoLowRhoCovA{}, at $\rho = 0.85$ from \AutoMidRhoCovH{} to \AutoMidRhoCovA{},
and at $\rho = 0.95$ from \AutoHighRhoCovH{} to \AutoHighRhoCovA{}. H-A2, that the
switch never overshoots nominal by more than five percent, is \AutoHtwoLabel{};
the worst coverage anywhere in the study is \AutoHtwoWorst{} against a ceiling of
\AutoHtwoCeiling{}. H-A3, that the improvement is larger at high branching ratio,
is \AutoHthreeLabel{}: the pooled difference is \AutoLowRhoDelta{} at $\rho =
0.3028$ and \AutoHighRhoDelta{} at $\rho = 0.95$.
